% ==============================================================================
% Fichier     : chapitres/08_reporting.tex
% Titre       : Rédaction du Rapport de Pentest et Remédiation
% Auteur      : MINKA MI NGUIDJOÏ Thierry Emmanuel

% ==============================================================================

\chapter{Rédaction du Rapport de Pentest et Remédiation}
\label{chap:reporting}

La phase technique est terminée. La \textbf{valeur ajoutée d'un pentester}
ne réside pas dans le nombre de machines compromises, mais dans la qualité
du rapport final. Ce document doit être compréhensible à la fois par les
équipes techniques (pour corriger) et par la direction (pour arbitrer les
budgets de sécurité). Un rapport de qualité professionnelle transforme des
\sigle{PoC} techniques en décisions d'investissement.

% ==============================================================================
\section{Structure du Rapport de Pentest}
\label{sec:rapport_structure}
% ==============================================================================

Le rapport suit une structure standardisée inspirée du \sigle{PTES}
(\textit{Reporting Phase}) et adaptée aux attentes des clients
professionnels. Il se compose de deux parties distinctes destinées à des
audiences différentes.

\subsection{Synthèse Exécutive (\textit{Executive Summary})}
\label{subsec:exec_summary}

La synthèse exécutive est destinée aux \textbf{décideurs} : \sigle{PDG},
\sigle{DSI}, \sigle{RSSI}, Conseil d'Administration. Elle doit tenir sur
1 à 2 pages maximum et ne contenir aucun jargon technique.

\begin{boxlizbiz}
\textbf{Template Executive Summary, LiZbiZ :}

\medskip
\textbf{Contexte et périmètre}

Du [DATE] au [DATE], le cabinet [NOM] a réalisé un test d'intrusion
de type \textit{Gray Box} sur l'infrastructure de LiZbiZ, conformément
au mandat signé le [DATE]. Le périmètre couvrait 8 systèmes incluant les
serveurs web, le portail employé, l'Active Directory et les extensions
cloud.

\medskip
\textbf{Résultats globaux}

\begin{center}
\begin{tabular}{lc}
\toprule
\textbf{Criticité} & \textbf{Nombre de vulnérabilités} \\
\midrule
Critique (CVSS 9.0--10.0) & 3 \\
Élevé (CVSS 7.0--8.9)     & 5 \\
Moyen (CVSS 4.0--6.9)     & 8 \\
Faible (CVSS 0.1--3.9)    & 4 \\
\midrule
\textbf{Total}             & \textbf{20} \\
\bottomrule
\end{tabular}
\end{center}

\medskip
\textbf{Constats majeurs}

L'infrastructure de LiZbiZ présente des vulnérabilités critiques permettant
à un attaquant externe d'obtenir le contrôle total du domaine Active Directory
en moins de 4~heures, sans exploiter de vulnérabilité de type \textit{zero-day}.
Les vecteurs principaux sont : un serveur web non mis à jour exposé sur Internet,
une vulnérabilité critique non patchée sur le service SMB Windows, et une
politique de mots de passe insuffisante sur les comptes de service.

\medskip
\textbf{Recommandations prioritaires}
\begin{enumerate}
    \item Appliquer immédiatement le patch \sigle{MS17-010} sur tous les
          serveurs Windows.
    \item Mettre à jour Apache~2.2.22 vers une version supportée (2.4.x+).
    \item Renforcer la politique de mots de passe des comptes de service
          (20 caractères minimum, rotation trimestrielle).
\end{enumerate}
\end{boxlizbiz}

\subsection{Méthodologie et Périmètre}
\label{subsec:methodologie_rapport}

Cette section détaille l'approche technique pour les équipes \sigle{IT} :
\begin{itemize}
    \item \textbf{Type de test :} Gray Box, compte \og{}stagiaire\fg{}
          fourni pour l'intranet.
    \item \textbf{Méthodologie :} PTES (\textit{Penetration Testing Execution
          Standard}), avec mapping \sigle{MITRE ATT\&CK} pour chaque
          technique employée.
    \item \textbf{Outils utilisés :} Nmap, Metasploit, SQLMap, BloodHound,
          Mimikatz, Chisel, GoPhish.
    \item \textbf{Périmètre testé :} Liste précise des \sigle{IP} et \sigle{URL}
          testées (référencer l'Annexe~A, tableau des cibles).
    \item \textbf{Périmètre exclu :} Actions DoS, COMEX, systèmes de paiement.
\end{itemize}

\subsection{Résultats Techniques (\textit{Findings})}
\label{subsec:findings}

C'est le cœur du rapport. Chaque vulnérabilité découverte fait l'objet
d'une fiche détaillée standardisée.

% ==============================================================================
\section{Le Score CVSS v3.1}
\label{sec:cvss}
% ==============================================================================

\begin{boxdef}
\textbf{CVSS v3.1} (\textit{Common Vulnerability Scoring System}, version~3.1,
publiée par le \sigle{FIRST} en 2019) : standard international de notation
de la gravité des vulnérabilités, de~0 (aucune) à~10 (critique). Le score
est calculé à partir de trois groupes de métriques : \textbf{Base},
\textbf{Temporal} et \textbf{Environmental}.
\end{boxdef}

\begin{table}[htbp]
\centering
\caption{Échelle de criticité CVSS v3.1}
\label{tab:cvss_scale}
\begin{tabular}{lll}
\toprule
\textbf{Score} & \textbf{Criticité} & \textbf{Action recommandée} \\
\midrule
0.0         & Informatif & Aucune urgence \\
0.1 -- 3.9  & Faible     & Traitement lors de la prochaine maintenance \\
4.0 -- 6.9  & Moyen      & Traitement sous 3 mois \\
7.0 -- 8.9  & Élevé      & Traitement sous 1 mois \\
9.0 -- 10.0 & Critique   & Traitement immédiat (sous 72h) \\
\bottomrule
\end{tabular}
\end{table}

\subsection{Métriques de Base CVSS v3.1}
\label{subsec:cvss_metrics}

\begin{table}[htbp]
\centering
\caption{Métriques Base CVSS v3.1 et leurs valeurs}
\label{tab:cvss_metrics}
\begin{tabular}{p{4cm}p{4cm}p{5.5cm}}
\toprule
\textbf{Métrique} & \textbf{Valeurs} & \textbf{Signification} \\
\midrule
Attack Vector (\sigle{AV})
    & Network / Adjacent / Local / Physical
    & Comment l'attaquant accède au système \\
\addlinespace
Attack Complexity (\sigle{AC})
    & Low / High
    & Conditions nécessaires à l'exploitation \\
\addlinespace
Privileges Required (\sigle{PR})
    & None / Low / High
    & Niveau de privilèges requis \\
\addlinespace
User Interaction (\sigle{UI})
    & None / Required
    & Interaction utilisateur nécessaire \\
\addlinespace
Scope (\sigle{S})
    & Unchanged / Changed
    & Impact au-delà du composant vulnérable \\
\addlinespace
Confidentiality (\sigle{C})
    & None / Low / High
    & Impact sur la confidentialité \\
\addlinespace
Integrity (\sigle{I})
    & None / Low / High
    & Impact sur l'intégrité \\
\addlinespace
Availability (\sigle{A})
    & None / Low / High
    & Impact sur la disponibilité \\
\bottomrule
\end{tabular}
\end{table}

\begin{boxinfo}
\textbf{CVSS v4.0, Évolution (novembre 2023) :}
La version~4.0 introduit quatre groupes de métriques : \textit{Base}
(inchangé), \textit{Threat} (exploitabilité réelle, remplace Temporal),
\textit{Environmental} (contexte spécifique de l'organisation) et
\textit{Supplemental} (informations additionnelles non scorées).
Elle distingue également les scores \og{}CVSS-B\fg{} (Base),
\og{}CVSS-BT\fg{} (Base+Threat) et \og{}CVSS-BE\fg{} (Base+Environmental).
Les rapports professionnels de 2024+ doivent préciser la version utilisée.
Ce cours utilise \textbf{CVSS v3.1} comme référence principale.
\end{boxinfo}

% ==============================================================================
\section{Modèle de Fiche de Vulnérabilité}
\label{sec:fiche_vuln}
% ==============================================================================

Chaque \textit{finding} doit suivre un format standardisé permettant une
lecture rapide par les équipes techniques et une priorisation claire par
la direction.

\begin{table}[htbp]
\centering
\caption{Fiche de constat technique, Injection SQL (LiZbiZ)}
\label{tab:fiche_sqli}
\begin{tabular}{p{3.5cm}p{10cm}}
\toprule
\textbf{Champ} & \textbf{Valeur} \\
\midrule
Titre
    & Injection SQL sur le formulaire de recherche \\
\addlinespace
Référence
    & FINDING-001 \\
\addlinespace
Sévérité
    & \textbf{CRITIQUE} \\
\addlinespace
Score CVSS v3.1
    & \textbf{9.8}, AV:N/AC:L/PR:N/UI:N/S:U/C:H/I:H/A:H \\
\addlinespace
CVE associée
    & N/A (vulnérabilité custom), CWE-89 (\textit{SQL Injection}) \\
\addlinespace
MITRE ATT\&CK
    & T1190 (\textit{Exploit Public-Facing Application}) \\
\addlinespace
Cible
    & \texttt{http://203.0.113.10/search.php} \\
\addlinespace
Description
    & Le paramètre \texttt{id} de la requête POST n'est pas filtré.
      Un attaquant peut injecter du code SQL arbitraire pour lire ou
      modifier la base de données sans authentification. \\
\addlinespace
Preuve (\sigle{PoC})
    & Capture d'écran (Annexe~B-001) montrant l'extraction de la table
      \texttt{users} via SQLMap. Hash SHA-256 de la preuve :
      \texttt{3a7f8d...} \\
\addlinespace
Impact business
    & Vol de l'intégralité des données clients (47 enregistrements
      dont données personnelles), exposition \sigle{RGPD}, amende
      potentielle jusqu'à 4\% du CA annuel. \\
\addlinespace
Recommandation
    & Utiliser des requêtes préparées (\textit{Prepared Statements})
      dans le code PHP. Implémenter un \sigle{WAF} en couche de
      compensation. Ne jamais concaténer les entrées utilisateur dans
      les requêtes SQL. \\
\addlinespace
Correction vs compensation
    & \textbf{Correction} : Prepared Statements (délai estimé : 2j).
      \textbf{Compensation} : Règle WAF bloquant les patterns SQLi. \\
\addlinespace
Référence
    & OWASP Top~10 2021, A03:2021 (Injection),
      OWASP Testing Guide v4.2 \S4.7.5. \\
\bottomrule
\end{tabular}
\end{table}

\begin{table}[htbp]
\centering
\caption{Fiche de constat technique, EternalBlue MS17-010}
\label{tab:fiche_eternal}
\begin{tabular}{p{3.5cm}p{10cm}}
\toprule
\textbf{Champ} & \textbf{Valeur} \\
\midrule
Titre
    & Absence de patch MS17-010 (EternalBlue) \\
\addlinespace
Référence
    & FINDING-002 \\
\addlinespace
Sévérité
    & \textbf{CRITIQUE} \\
\addlinespace
Score CVSS v3.1
    & \textbf{9.8}, AV:N/AC:L/PR:N/UI:N/S:U/C:H/I:H/A:H \\
\addlinespace
CVE associée
    & CVE-2017-0144 \\
\addlinespace
MITRE ATT\&CK
    & T1210 (\textit{Exploitation of Remote Services}),
      T1569.002 (\textit{Service Execution}) \\
\addlinespace
Cible
    & 192.168.10.20 (Windows Server 2012 R2) \\
\addlinespace
Description
    & Vulnérabilité de débordement de tampon dans SMBv1 permettant
      l'exécution de code arbitraire à distance sans authentification.
      Patch disponible depuis mars 2017 (MS17-010). \\
\addlinespace
Preuve (\sigle{PoC})
    & Capture d'écran (Annexe~B-002) : shell SYSTEM obtenu en 28
      secondes. Output de \texttt{getuid} : \texttt{NT AUTHORITY\textbackslash{}SYSTEM}.
      Hash SHA-256 : \texttt{8b2c4f...} \\
\addlinespace
Impact business
    & Contrôle total du serveur Windows, accès aux données RH,
      possibilité de ransomware, point de pivot vers l'AD. \\
\addlinespace
Recommandation
    & Appliquer immédiatement le patch MS17-010. Désactiver SMBv1
      sur tous les systèmes (sauf nécessité opérationnelle documentée).
      Segmenter le réseau pour isoler les serveurs Windows. \\
\addlinespace
Correction vs compensation
    & \textbf{Correction} : Patch MS17-010 (délai : 24h).
      \textbf{Compensation} : Bloquer le port 445 au niveau du
      pare-feu interne. \\
\bottomrule
\end{tabular}
\end{table}

% ==============================================================================
\section{Plan de Remédiation et Priorisation}
\label{sec:remediation}
% ==============================================================================

\subsection{Matrice de Priorisation}
\label{subsec:priorisation}

Le plan de remédiation doit être actionnable et priorisé selon le risque
réel, non selon la complexité technique de la correction.

\begin{table}[htbp]
\centering
\caption{Plan de remédiation LiZbiZ, priorisé par criticité et effort}
\label{tab:plan_remediation}
\begin{tabular}{p{0.4cm}p{3.5cm}p{1.5cm}p{2cm}p{4.5cm}}
\toprule
\textbf{\#} & \textbf{Vulnérabilité} & \textbf{CVSS} & \textbf{Délai} & \textbf{Action corrective} \\
\midrule
1 & MS17-010 EternalBlue
    & 9.8 Critique
    & 24h
    & Patch MS17-010, désactiver SMBv1 \\
\addlinespace
2 & Injection SQL / login
    & 9.8 Critique
    & 2j
    & Prepared Statements + WAF \\
\addlinespace
3 & Apache 2.2.22 (RCE)
    & 9.8 Critique
    & 48h
    & Mise à jour vers Apache 2.4.x \\
\addlinespace
4 & Phishing (taux 34\%)
    & 8.1 Élevé
    & 1 mois
    & Sensibilisation + MFA \\
\addlinespace
5 & LDAP anonymous bind
    & 7.5 Élevé
    & 1 mois
    & Désactiver l'accès anonyme LDAP \\
\addlinespace
6 & Politique mdp comptes service
    & 7.2 Élevé
    & 1 mois
    & 20 chars min + rotation 90j \\
\addlinespace
7 & Répertoire /backup/ exposé
    & 6.5 Moyen
    & 3 mois
    & Restriction d'accès + .htaccess \\
\addlinespace
8 & XSS stored intranet
    & 6.1 Moyen
    & 3 mois
    & Encodage sorties HTML + CSP \\
\bottomrule
\end{tabular}
\end{table}

\subsection{Correction vs Compensation}
\label{subsec:correction_compensation}

\begin{boxdef}
\textbf{Correction :} Suppression de la faille à la source, mise à jour
du logiciel, correction du code, reconfiguration.

\textbf{Compensation (Mitigation) :} Barrière ajoutée en couche de
protection lorsque la correction directe est impossible ou trop longue
— règle de pare-feu, \sigle{WAF}, désactivation du service.
La compensation ne supprime pas la vulnérabilité : elle réduit la
surface d'exposition dans l'attente de la correction définitive.
\end{boxdef}

% ==============================================================================
\section{Livrable Final et Débriefing}
\label{sec:livrable_debrief}
% ==============================================================================

\subsection{Structure des Livrables}
\label{subsec:livrables_finaux}

\begin{table}[htbp]
\centering
\caption{Livrables finaux de la mission LiZbiZ}
\label{tab:livrables_finaux}
\begin{tabular}{p{3.5cm}p{2.5cm}p{7cm}}
\toprule
\textbf{Livrable} & \textbf{Format} & \textbf{Contenu} \\
\midrule
Rapport Exécutif
    & PDF (5--10 pages)
    & Synthèse, tableau de criticité, 3 recommandations prioritaires \\
\addlinespace
Rapport Technique
    & PDF (20--40 pages)
    & Méthodologie, 20 fiches de vulnérabilités, mapping MITRE,
      plan de remédiation priorisé \\
\addlinespace
Annexe A, Périmètre
    & PDF
    & Liste complète des cibles testées avec adresses IP et URLs \\
\addlinespace
Annexe B, Preuves
    & ZIP chiffré
    & Captures d'écran horodatées + hashes \sigle{SHA-256},
      fichiers flags, logs de sessions \\
\addlinespace
Scripts \sigle{PoC}
    & ZIP chiffré
    & Exploits utilisés, scripts personnalisés, instructions de
      reproduction \\
\addlinespace
Présentation Débriefing
    & \sigle{PPTX}/\sigle{PDF}
    & Support 15 slides pour la réunion de restitution \\
\bottomrule
\end{tabular}
\end{table}

\begin{boxinfo}
\textbf{Chaîne de custody des livrables :} Chaque fichier de preuve
doit être accompagné de son hash \sigle{SHA-256} calculé immédiatement
après capture (Section~\ref{sec:preuves}). L'archive ZIP contenant les
preuves doit être chiffrée (\sigle{AES-256}) et transmise de façon
sécurisée (canal chiffré, jamais par email non chiffré).

\begin{lstlisting}[language=bash]
# Créer l'archive chiffrée des preuves
zip -e -r preuves_lizbiz_YYYYMMDD.zip ./preuves/
# Ou avec gpg (symétrique)
tar czf - ./preuves/ | gpg -c --cipher-algo AES256 \
    -o preuves_lizbiz_YYYYMMDD.tar.gz.gpg

# Calculer le hash de l'archive finale
sha256sum preuves_lizbiz_YYYYMMDD.tar.gz.gpg >> livraison_manifest.txt
\end{lstlisting}
\end{boxinfo}

\subsection{La Réunion de Débriefing}
\label{subsec:debrief}

La mission se conclut par une \textbf{réunion de restitution} structurée
en trois temps :

\begin{enumerate}
    \item \textbf{Présentation exécutive} (30~min) : synthèse des risques
          majeurs pour la direction, sans jargon technique. Utiliser
          des analogies métier pour expliquer l'impact (ex.~:
          \og{}EternalBlue permettrait à un ransomware de chiffrer tous
          vos serveurs Windows en moins de 10~minutes\fg{}).
    \item \textbf{Revue technique} (60~min) : démonstration des \sigle{PoC}
          avec les équipes \sigle{IT}/\sigle{SOC}. Chaque finding est
          présenté avec sa preuve, son score \sigle{CVSS} et la
          contre-mesure recommandée.
    \item \textbf{Plan d'action} (30~min) : validation collective du
          planning de remédiation et désignation des propriétaires de
          chaque action corrective.
\end{enumerate}

\begin{boxalert}
\textbf{Transmission sécurisée du rapport :} Le rapport de pentest contient
des informations hautement sensibles sur les vulnérabilités du client.
Il ne doit \textbf{jamais} être transmis par email non chiffré. Utiliser
un canal chiffré \sigle{E2E} ou une plateforme de partage sécurisée
(\sigle{SFTP}, portail chiffré). Le rapport doit être supprimé des
systèmes du prestataire à l'issue de la mission selon les termes du
contrat de confidentialité.
\end{boxalert}

% ==============================================================================
\section*{Points Clés à Retenir}
% ==============================================================================

\begin{boxinfo}
\textbf{Ce qu'il faut retenir de ce chapitre :}
\begin{enumerate}
    \item Le rapport est le \textbf{produit final} du pentest, un
          excellent rapport sur des résultats modestes vaut mieux qu'un
          mauvais rapport sur des résultats brillants.
    \item L'\textbf{Executive Summary} ne doit contenir aucun jargon
          technique, il s'adresse aux décideurs, pas aux ingénieurs.
    \item Chaque fiche de vulnérabilité doit inclure \textbf{CVE/CWE,
          CVSS v3.1 avec vecteur complet, référence MITRE ATT\&CK,
          PoC avec hash SHA-256, impact business, et recommandation}.
    \item La distinction \textbf{Correction vs Compensation} est
          fondamentale, elle donne de la flexibilité au client sur
          le délai de remédiation.
    \item Les livrables doivent être \textbf{chiffrés et signés} avant
          transmission, le rapport expose toutes les failles du client.
    \item \textbf{CVSS v3.1} est la version de référence du cours ;
          CVSS v4.0 (nov.~2023) sera progressivement adopté, toujours
          préciser la version dans le rapport.
\end{enumerate}
\end{boxinfo}

% ==============================================================================
\section*{Exercice Pratique}
% ==============================================================================

\begin{boxexo}
\textbf{Exercice 8.1, Rédaction d'un rapport de pentest LiZbiZ}

\medskip
\textbf{Objectif :} Produire un rapport de pentest professionnel complet
sur la base des résultats des exercices précédents.

\medskip
\textbf{Tâches :}
\begin{enumerate}
    \item Rédiger une \textbf{synthèse exécutive} d'une page (tableau de
          criticité + 3 recommandations prioritaires en langage non technique).
    \item Compléter \textbf{5 fiches de vulnérabilités} au format
          Tableau~\ref{tab:fiche_sqli}, incluant :
          CVE/CWE, score \sigle{CVSS v3.1} avec vecteur complet,
          référence MITRE ATT\&CK, capture d'écran \sigle{PoC} avec
          hash \sigle{SHA-256}, impact business, recommandation
          correction + compensation.
    \item Établir le \textbf{plan de remédiation priorisé} au format
          Tableau~\ref{tab:plan_remediation}.
    \item Créer l'archive chiffrée des preuves et calculer son
          hash SHA-256.
\end{enumerate}

\medskip
\textbf{Livrable :} Rapport PDF complet (10--15 pages) + archive ZIP
chiffrée des preuves. Le rapport sera évalué sur la clarté de
l'Executive Summary, la complétude des fiches et la précision des
scores CVSS.
\end{boxexo}
